\chapter*{ВВЕДЕНИЕ}
\addcontentsline{toc}{chapter}{ВВЕДЕНИЕ}

Plankalkul является одним из первых языков программирования высокого уровня.
Язык был предложен Конрадом Цузе для описания алгоритмов независимо от конкретного устройства вычислительной машины.
Несмотря на исторический характер языка, его идеи остаются полезными для изучения ранних подходов к типизации, описанию программ, работе с составными данными и разделению алгоритма на подпрограммы.

Исходная запись Plankalkul была рассчитана на двумерную математическую нотацию.
Для практической реализации компилятора требуется линейное текстовое представление, пригодное для обработки современными средствами синтаксического анализа.
В данной работе рассматривается такая версия языка: она сохраняет основные классы переменных Plankalkul, типы на основе битов, массивы, записи, условные операторы и циклы, но описывается в виде ANTLR4-грамматики.

Целевой платформой компиляции выбран Common Intermediate Language платформы .NET.
Такой выбор позволяет представить каждую программу Plankalkul в виде статического метода CLR, использовать стандартные типы времени выполнения и выполнять сгенерированный код после сборки средствами .NET.

Цель работы -- разработать компилятор языка Plankalkul в CIL.

Для достижения поставленной цели необходимо решить следующие задачи:
\begin{enumerate}
	\item описать грамматику языка Plankalkul, целевую платформу CIL и сформулировать требования к компилятору;
	\item спроектировать компилятор, включающую фронтенд, промежуточное представление, семантический анализ и генерацию CIL;
	\item реализовать обработку программ на языке Plankalkul и интерфейс командной строки;
	\item проверить корректность синтаксического и семантического анализа на тестовых программах.
\end{enumerate}
